% GENERATED FILE. Edit canonical lessons, then run npm run generate:books.
% canonical-source-hash: fnv1a64:4f999e5a198d2fb3
% canonical-lessons: AR-C208-mana, AR-C208-mithal, AR-C208-hurriyya, AR-C208-qaida, AR-C208-tariqa

\chapter{Meaning, Example, Freedom, Rule, Way}
\label{ch:ar-meaning-example-freedom-rule-way}
\hlchaptermodality{208}

\begin{chapteropening}
\textbf{By the end of this chapter:} \emph{I can say a meaning, an example, freedom, a rule and a way.}

Complete the last lesson of chapter 208: I can say a meaning, an example, freedom, a rule and a way.
\end{chapteropening}

\section[\texorpdfstring{\ar{معنى} (maʿnā)}{معنى (maʿnā)}]{\ar{معنى} (maʿnā) — a meaning}

\label{lesson:AR-C208-mana}



\begin{quote}
Before the new one: say the Arabic for if, then the Arabic for as.
\end{quote}

\subsection*{You'll want to know: \ar{معنى}}
\textbf{\ar{معنى}} — \emph{maʿnā} — \textquotedblleft{}a meaning\textquotedblright{}.

\begin{cousinweb}
\textbf{\ar{معنى}} (\emph{maʿnā}) is built on the native Arabic root \emph{ʿ-n-y}, \textquotedblleft{}to mean, to intend\textquotedblright{} (Lane). \emph{Maʿnā} is \textquotedblleft{}what is meant\textquotedblright{}.
\end{cousinweb}

\begin{grammarlens}[title={What you've built}]
One more word for this chapter.
\end{grammarlens}

\subsection*{Your turn}
\begin{itemize}
\raggedright
  \item \emph{Say it:} \emph{maʿnā}
  \item \emph{Say it:} \emph{maʿnā}, once more
  \item \emph{Say it:} say \emph{maʿnā}
  \item \emph{Recall:} say the Arabic for if, then the Arabic for as, then say \emph{maʿnā} again
\end{itemize}

\begin{tcolorbox}[breakable,colback=teal!4,colframe=teal!35!black,title={Before you move on}]
What does \ar{معنى} mean? (\textquotedblleft{}A meaning\textquotedblright{}.) Say it once more.
\end{tcolorbox}
\section[\texorpdfstring{\ar{مثال} (mithāl)}{مثال (mithāl)}]{\ar{مثال} (mithāl) — an example}

\label{lesson:AR-C208-mithal}



\begin{quote}
Before the new one: say the Arabic for as, then the Arabic for a meaning.
\end{quote}

\subsection*{You'll want to know: \ar{مثال}}
\textbf{\ar{مثال}} — \emph{mithāl} — \textquotedblleft{}an example\textquotedblright{}.

\begin{cousinweb}
\textbf{\ar{مثال}} (\emph{mithāl}) is built on the native Arabic root \emph{m-th-l}, \textquotedblleft{}to be like\textquotedblright{} (Lane). \emph{Mithāl} is a likeness, a pattern, hence \textquotedblleft{}an example\textquotedblright{}.
\end{cousinweb}

\begin{grammarlens}[title={What you've built}]
One more word for this chapter.
\end{grammarlens}

\subsection*{Your turn}
\begin{itemize}
\raggedright
  \item \emph{Say it:} \emph{mithāl}
  \item \emph{Say it:} \emph{mithāl}, once more
  \item \emph{Say it:} say \emph{mithāl}
  \item \emph{Recall:} say the Arabic for as, then the Arabic for a meaning, then say \emph{mithāl} again
\end{itemize}

\begin{tcolorbox}[breakable,colback=teal!4,colframe=teal!35!black,title={Before you move on}]
What does \ar{مثال} mean? (\textquotedblleft{}An example\textquotedblright{}.) Say it once more.
\end{tcolorbox}
\section[\texorpdfstring{\ar{حرية} (ḥurriyya)}{حرية (ḥurriyya)}]{\ar{حرية} (ḥurriyya) — freedom}

\label{lesson:AR-C208-hurriyya}



\begin{quote}
Before the new one: say the Arabic for a meaning, then the Arabic for an example.
\end{quote}

\subsection*{You'll want to know: \ar{حرية}}
\textbf{\ar{حرية}} — \emph{ḥurriyya} — \textquotedblleft{}freedom\textquotedblright{}.

\begin{cousinweb}
\textbf{\ar{حرية}} (\emph{ḥurriyya}) is built on the native Arabic root \emph{ḥ-r-r}, \textquotedblleft{}to be free; to be hot\textquotedblright{} (Lane). \emph{Ḥurriyya} is \emph{ḥurr}, \textquotedblleft{}free\textquotedblright{}, with the abstract ending \emph{-iyya}.
\end{cousinweb}

\begin{grammarlens}[title={What you've built}]
One more word for this chapter.
\end{grammarlens}

\subsection*{Your turn}
\begin{itemize}
\raggedright
  \item \emph{Say it:} \emph{ḥurriyya}
  \item \emph{Say it:} \emph{ḥurriyya}, once more
  \item \emph{Say it:} say \emph{ḥurriyya}
  \item \emph{Recall:} say the Arabic for a meaning, then the Arabic for an example, then say \emph{ḥurriyya} again
\end{itemize}

\begin{tcolorbox}[breakable,colback=teal!4,colframe=teal!35!black,title={Before you move on}]
What does \ar{حرية} mean? (\textquotedblleft{}Freedom\textquotedblright{}.) Say it once more.
\end{tcolorbox}
\section[\texorpdfstring{\ar{قاعدة} (qāʿida)}{قاعدة (qāʿida)}]{\ar{قاعدة} (qāʿida) — a rule}

\label{lesson:AR-C208-qaida}



\begin{quote}
Before the new one: say the Arabic for an example, then the Arabic for freedom.
\end{quote}

\subsection*{You'll want to know: \ar{قاعدة}}
\textbf{\ar{قاعدة}} — \emph{qāʿida} — \textquotedblleft{}a rule\textquotedblright{}.

\begin{cousinweb}
\textbf{\ar{قاعدة}} (\emph{qāʿida}) is built on the native Arabic root \emph{q-ʿ-d}, \textquotedblleft{}to sit\textquotedblright{} (Lane). Lane's \emph{qawāʿid} are the foundations a house sits on; \emph{qāʿida} is a base, a rule.
\end{cousinweb}

\begin{grammarlens}[title={What you've built}]
One more word for this chapter.
\end{grammarlens}

\subsection*{Your turn}
\begin{itemize}
\raggedright
  \item \emph{Say it:} \emph{qāʿida}
  \item \emph{Say it:} \emph{qāʿida}, once more
  \item \emph{Say it:} say \emph{qāʿida}
  \item \emph{Recall:} say the Arabic for an example, then the Arabic for freedom, then say \emph{qāʿida} again
\end{itemize}

\begin{tcolorbox}[breakable,colback=teal!4,colframe=teal!35!black,title={Before you move on}]
What does \ar{قاعدة} mean? (\textquotedblleft{}A rule\textquotedblright{}.) Say it once more.
\end{tcolorbox}
\section[\texorpdfstring{\ar{طريقة} (ṭarīqa)}{طريقة (ṭarīqa)}]{\ar{طريقة} (ṭarīqa) — a way, a method}

\label{lesson:AR-C208-tariqa}



\begin{quote}
Before the new one: say the Arabic for freedom, then the Arabic for a rule.
\end{quote}

\subsection*{You'll want to know: \ar{طريقة}}
\textbf{\ar{طريقة}} — \emph{ṭarīqa} — \textquotedblleft{}a way, a method\textquotedblright{}.

\begin{cousinweb}
\textbf{\ar{طريقة}} (\emph{ṭarīqa}) is built on the native Arabic root \emph{ṭ-r-q}, \textquotedblleft{}to strike, to beat\textquotedblright{} (Lane). Lane's \emph{ṭarīqa} is a way, a manner of conduct.
\end{cousinweb}

\begin{grammarlens}[title={What you've built}]
One more word for this chapter.
\end{grammarlens}

\subsection*{Your turn}
\begin{itemize}
\raggedright
  \item \emph{Say it:} \emph{ṭarīqa}
  \item \emph{Say it:} \emph{ṭarīqa}, once more
  \item \emph{Say it:} say \emph{ṭarīqa}
  \item \emph{Recall:} say the Arabic for freedom, then the Arabic for a rule, then say \emph{ṭarīqa} again
\end{itemize}

\begin{tcolorbox}[breakable,colback=teal!4,colframe=teal!35!black,title={Before you move on}]
What does \ar{طريقة} mean? (\textquotedblleft{}A way, a method\textquotedblright{}.) Say it once more.
\end{tcolorbox}
